\documentclass[11pt]{article}
\usepackage[margin=1in]{geometry}
\usepackage{amsmath,amssymb,amsthm}
\usepackage{booktabs}

\newtheorem*{claim}{Conjecture (00000000477)}

\newcommand{\syt}[5]{$\begin{array}{@{}ccc@{}}#1&#2&#3\\#4&#5&\end{array}$}

\title{Disproof of TLMC Conjecture 00000000477\\
\large Promotion orbit lengths on linear extensions}
\author{}
\date{}

\begin{document}
\maketitle

\section{Verdict}

\begin{center}
\fbox{\textbf{FALSE.} The lcm of promotion orbit lengths need not divide $\#\mathrm{LE}(P)$.}
\end{center}

\begin{claim}
Let $P$ be a finite poset and let promotion act on the set $\mathrm{LE}(P)$ of its
linear extensions.  Then the lcm of the promotion orbit lengths divides
$\#\mathrm{LE}(P)$.
\end{claim}

\section{The counterexample}

Take $P$ to be the Ferrers poset of the Young diagram of shape $(3,2)$: its five
elements are the cells
$(0,0),(0,1),(0,2),(1,0),(1,1)$ of the diagram, ordered componentwise, i.e.\
$(i,j) \le (i',j')$ iff $i \le i'$ and $j \le j'$.

Linear extensions of $P$ are in bijection with standard Young tableaux (SYT) of
shape $(3,2)$: the entry of a cell is its position in the extension.  The hook
length formula gives
\[
\#\mathrm{LE}(P) \;=\; \frac{5!}{4\cdot 3\cdot 1 \cdot 2 \cdot 1} \;=\; 5,
\]
and indeed exactly five tableaux exist:
\begin{center}
\begin{tabular}{ccccc}
\toprule
$t_1$ & $t_2$ & $t_3$ & $t_4$ & $t_5$\\
\midrule
\syt{1}{2}{3}{4}{5} &
\syt{1}{2}{4}{3}{5} &
\syt{1}{2}{5}{3}{4} &
\syt{1}{3}{4}{2}{5} &
\syt{1}{3}{5}{2}{4} \\\bottomrule
\end{tabular}
\end{center}

\section{Promotion and its orbits}

Apply Sch\"utzenberger promotion (remove the corner entry, jeu-de-taquin slide
the hole to an outer corner, write $6$ in the hole, decrease all entries by $1$).
Working out the jeu de taquin for this fixed shape (the hole makes at most two
slides), one obtains
\[
t_1 \mapsto t_3 \mapsto t_4 \mapsto t_1,
\qquad
t_2 \mapsto t_5 \mapsto t_2 .
\]
For instance, promoting $t_1 = (1,2,3,4,5)$: the right neighbour $2$ of the hole
is smaller than the lower neighbour $4$, so $2$ slides up; then of $3$ (right)
and $5$ (below), the smaller $3$ slides; the hole reaches the outer corner
$(0,2)$, receives $6$, and decreasing by $1$ yields $(1,2,5,3,4) = t_3$.

So promotion partitions the five linear extensions into two orbits, of lengths
$3$ and $2$:
\[
\{t_1, t_3, t_4\}, \qquad \{t_2, t_5\}.
\]

\section{The falsification}

\[
\operatorname{lcm}(3,2) = 6, \qquad 6 \nmid 5 = \#\mathrm{LE}(P).
\]
Hence the conjecture fails already for this five-element poset.

\section{Verification}

\begin{itemize}
  \item \texttt{reproduce.py}: an independent brute-force recomputation.  It
  enumerates all $5!$ orderings of the cells, keeps the linear extensions
  ($5$ of them), converts them to tableaux, applies Sch\"utzenberger promotion,
  and recovers the orbit lengths $[2,3]$, lcm $6$, and $6 \nmid 5$.
  Run with \texttt{python3 reproduce.py}; assertions check every number.
  \item \texttt{lean4/}: a core-Lean-4 (no Mathlib) certificate.  \texttt{Main.lean}
  enumerates the $5^5$ candidate tuples in the kernel and proves
  \texttt{syts = [t1,\dots,t5]} and \texttt{syts.length = 5}; promotion is
  implemented explicitly and the orbit equations
  $t_1 \to t_3 \to t_4 \to t_1$, $t_2 \to t_5 \to t_2$ plus
  $\operatorname{lcm}(3,2)=6$ and $6 \nmid 5$ are proved by \texttt{decide}.
  \texttt{Check.lean} audits every theorem with \texttt{\#print axioms}: all
  report \emph{does not depend on any axioms} (no \texttt{propext}, no
  \texttt{Quot.sound}, no \texttt{Classical.choice}), no \texttt{sorry}.
\end{itemize}

\section{Scope and boundary}

The conjecture is stated for arbitrary finite posets, so the single counterexample
above refutes it.  It is worth recording where the divisibility does hold: for
rectangular shapes $(a^b)$ promotion is known (Reiner--Stanton--White cyclic
sieving; Haiman) to have a single orbit of length $ab$, whose lcm $ab$ trivially
divides the number of SYT only in special cases --- the point here is merely that
the non-rectangular shape $(3,2)$ breaks the claimed divisibility, with orbit
lengths $3+2=5=\#\mathrm{LE}(P)$ whose lcm $6$ exceeds $\#\mathrm{LE}(P)$ itself.

\end{document}
